\section{问题重述与分析}

本题要求建立从多源遥感识别到灾害评估和救援决策的完整链条。前两问以灾前光学、灾后SAR和新增淹没标签为基础，重点解决跨传感器融合、类别不平衡、空间泛化和区域面积误差。第三问将模型迁移到2026年7月6日前后发生洪水的广西横州市，需要独立选择研究区和事件前后影像，区分既有水体与新增淹没，并计算新增面积、消退面积、净增加量和面积比例。

问题二和问题三的关键难点有三点：第一，训练数据与横州应用数据来自不同传感器，存在显著域偏移；第二，面积计算不能直接使用经纬度像素差，必须统一到冻结的投影参考栅格；第三，洪灾强度既要同时反映绝对淹没面积和局地覆盖比例，又要将识别误差传播到等级置信度，避免把接近分级阈值的网格解释为确定结论。
